\section{Вопросы без ответов}

\subsection{Чем \(d\) отличается от \(p\)?}
Оба уровня --- и \(d\), и \(p\) --- лежат за пределами всего счётного. Спросите Вселенную,
спросите свои чувства, спросите почти всю математику: об обоих вы услышите один ответ ---
«слишком много». Разницы не почувствовать. И всё же \(p\) чудовищно больше \(d\).
Насколько? Чтобы ужать \(p\) до размеров \(d\), пришлось бы взять от него логарифм
\emph{двенадцать} раз подряд. А \(d\) --- это уже число, у которого число цифр в числе его
цифр астрономично. При этом \(p+d=p\) и \(p-d=p\) с невообразимой точностью: в арифметике
этих гигантов \(d\) --- погрешность округления. И одновременно \(d\) больше всего, что можно
построить из атомов Вселенной. Как удержать обе мысли разом?

\subsection{Как с этим работать?}
Никак --- если «работать» значит «записать» или «вычислить»: не хватит ни бумаги, ни
атомов, ни времени. Математик работает не с числом, а с его \emph{адресом}: высотой башни,
скоростью роста, ординалами как линейкой. Мы сравниваем не числа, а темпы их бегства в
бесконечность. Но возиться с адресом числа --- то же ли, что знать число? Мы уверенно пишем
\(\ccoma\), ни разу его не увидев. Что у нас в руках --- число или имя?

\subsection{Существует ли оно?}
\(d\) нельзя записать, нельзя досчитать до него, нельзя хранить. Считай Вселенная от единицы
вечно --- она бы до \(d\) не дошла. В каком смысле \(d\) существует? Мы его придумали --- или
обнаружили? Если это факт о мире, то где он лежит? Если только о нашей записи --- почему он
подчиняется тем же законам, что и настоящие числа?

\subsection{Список вопросов}
\begin{itemize}
  \item Если и \(d\), и \(p\) для человека «бесконечно большие», есть ли между ними
        разница --- или она живёт только на бумаге?
  \item Число, которое нельзя ни записать, ни вычислить, ни досчитать, --- существует ли оно?
  \item Мы придумали крипто кому --- или наткнулись на неё?
  \item Работать с адресом числа --- это работать с числом?
  \item «Мало» --- мало для чего? Что добавляет размер, чего не хватает смыслу?
  \item Есть ли последнее число --- или у мысли нет потолка?
  \item Почему нас вообще тянет вверх по этой лестнице?
\end{itemize}
